%% ===========================================================================
\FloatBarrier
\section{Conclusions}
\label{sec:conclusions}
%% ===========================================================================

The formulation combines piezoelectric and flexoelectric coupling with
fractional memory in a $C^1$ Bogner--Fox--Schmit discretization. The second
moment of the averaging window sets the gradient stiffness, while modal L1
integration retains the dielectric history. The principal findings are:

\begin{itemize}
\item \textbf{The averaging window, through its second-moment tensor $M$,
  supplies the internal length.} Physically, the window is the domain over which
  the microstructure is averaged, and because the gradient energy is linear in
  $M$ the shape of that domain is not a parameter: windows of equal second
  moment give identical solutions to the third decimal, while their size spans
  the coupling from $3.6\times10^{-6}$ to $1.0\times10^{-2}$ over the map
  computed here. Rotating an anisotropic window changes the extensional
  coupling by $15$--$19\,\%$ between the aligned and transverse positions,
  because the gradient stiffness samples different directional strain gradients.
\item \textbf{The thickness $h$} converts a curvature into a strain gradient
  of order $\kappa_c/h$ and thereby separates the two couplings: the
  flexoelectric bending coupling grows by a factor $24$ between
  \SI{10}{\micro\meter} and \SI{1}{\micro\meter}, while the piezoelectric
  extensional coupling stays constant to $0.09\,\%$ and its quasi-static value,
  $0.5581$, lies within $0.5\,\%$ of the closed form $0.5608$.
\item \textbf{The flexoelectric coefficient $\mu$}, the polarization produced by
  a unit strain gradient, fixes the absolute size of the bending coupling: at
  the assumed $\mu=10^{-6}$\,\si{\coulomb\per\meter} the ratio of the two
  couplings spans $3.1\times10^{-4}$ to $7.5\times10^{-3}$ and reaches
  $0.1\,\%$ only below $h=5.5\,\si{\micro\meter}$.
\item \textbf{The Caputo order $\alpha$}, which measures how broad the
  relaxation spectrum is, controls the shape of the decay: the first-period
  decrement of the tip displacement runs from $0.063$ at $\alpha=1$ to $0.104$
  at $\alpha=0.6$ and back to $0.085$ at $\alpha=0.4$, so a free decay is not
  exponential below $\alpha=1$ and an integer-order reading is in error by
  $64\,\%$ at $\alpha=0.6$.
\item \textbf{The relaxation time $\tau$}, which places the loss peak at
  $\omega\tau$ of order unity, fixes where that damping acts: at $\alpha=1$ and
  $\omega\tau=0.3$ the damping ratio is $0.15003$ on the condensed matrix
  against $0.0675$ on the short-circuit one, a ratio of $0.450$, and the
  Kelvin--Voigt limit is reproduced to $0.02\,\%$.
\item \textbf{The element that carries all of it}, a $C^1$ Bogner--Fox--Schmit
  rectangle, converges at the $N^{-2}$ rate of bicubic Hermite interpolation, so
  that $48\times4$ elements reproduce the coupling factor to
  $6.2\times10^{-4}$ relative.
\end{itemize}

For sensing design, select the electrodes first ($0.4580$ against $0.1141$
at equal tip motion), then the thickness and window geometry; interpret decay
with $\alpha$, $\tau$ and the memory matrix stated together. Applications
remain conditional on the assumed inputs and two-dimensional kinematics:
plane-stress coupling differs from the three-dimensional value by $2.8\,\%$,
motivating measured-window calibration and shell extensions.

\FloatBarrier
